\chapter{SPARK simulation: algorithm, results and verification}
\label{sec:hf-spark-study}
\section{Source context, cases and random-state construction}
The native evidence used here is the final Hoeflin run of the 4 November
2026 five-contract campaign. In repository-relative notation, its directory is
\begin{center}
\small\path{verification/evidence/five-contracts-computer10-20261104/}\par
\path{runtime-final/hoeflin/check/}
\end{center}
The relevant files are \path{simulation.csv}, \path{calculations.csv},
\path{runtime.json}, \path{validation-normal.json} and
\path{replay-normal.json}, with corresponding optimized-validator records.
The sibling \path{ignore/} directory records the Ghost Ignore execution.
The campaign's \path{manifest.json}, native portfolio result and
\path{runtime-final/summary.json} identify the source hashes, tools,
commands and policy comparison. The simulation CSV SHA-256 is
\begin{center}\small
\path{dba8b3046072a6f46b4c63117df9ac0b33aaa4f486507c493558391395c0ec36}.
\end{center}
It is byte-identical to the earlier archived CSV used as the Dask
campaign's historical input. No retained evidence file is rewritten by
this report revision, and no new native execution is claimed.

The SPARK driver evaluates all eight specified battery sizes at
\(c=0.6\), using \(T=1{,}000{,}000\) draws for each size. For every
size it resets a 64-bit state to seed \(20261030\). The recorded generator
advances by the modular recurrence
\begin{align*}
 s_{j+1}&=(A s_j+B)\bmod 2^{64},\\
 A&=6364136223846793005,\qquad B=1442695040888963407.
\end{align*}
It maps the high 52 bits into an open-unit value
\[
 U_j=\bigl(\lfloor s_j/2^{12}\rfloor+\tfrac12\bigr)2^{-52}.
\]
Because \(0\leq\lfloor s_j/2^{12}\rfloor\leq2^{52}-1\),
\(2^{-53}\leq U_j\leq1-2^{-53}\). This elementary bound explains
why the subsequent logarithm avoids zero; it does not establish a
continuous uniform law or independent successive values.

Two successive values \(U,V\) give the Box--Muller pair
\[
 G=\sqrt{-2\log U}\cos(2\pi V),\qquad
 E=\sqrt{-2\log U}\sin(2\pi V).
\]
The driver then computes \(M_n\) by \eqref{eq:hf-aggregate}, divides by
\(\sqrt{v_n}\), accumulates its first two sample moments, and updates the
three nested threshold counts. Resetting the seed reuses the same base
pairs across sizes. Consequently the SPARK experiment shares the same
fixed-size aggregate coupling described for Dask, rather than simulating
a nested sequence of independent individual-test residuals.
Proposition~\ref{hfpr:simulation} proves the coupling distinction.

\begin{proposition}[Ideal Box--Muller law]\label{hfsp:box-muller}
If \(U,V\) are independent continuous uniform variables on \((0,1)\),
the displayed Box--Muller formulas yield independent standard normals.
\end{proposition}
\begin{proof}
This is the classical transformation of \citet{boxmuller1958hf}; the
following change-of-variables argument gives the complete law. Put
\(R=\sqrt{-2\log U}\) and \(\Theta=2\pi V\). Since
\(U=\exp(-R^2/2)\), differentiation gives radial density
\(r\exp(-r^2/2)\) for \(r>0\). The angle has density \(1/(2\pi)\)
on \((0,2\pi)\), independently of \(R\). The polar map
\((r,\theta)\mapsto(g,e)=(r\cos\theta,r\sin\theta)\)
has absolute Jacobian \(r\), outside a measure-zero origin and ray.
The density change-of-variables theorem therefore gives joint density
\[
 \frac{1}{2\pi}\exp\!\left(-\frac{g^2+e^2}{2}\right)
 =\phi(g)\phi(e).
\]
Factorization into the two normalized marginal densities proves both
standard-normality and independence.
\end{proof}
The proposition is conditional on ideal continuous inputs. The finite
LCG, discrete open-unit mapping and external elementary functions are not
proved to satisfy those hypotheses. Indeed, the displayed lower bound
on \(U\) gives \(R\leq\sqrt{106\log2}\) in exact arithmetic and hence
\(|G|,|E|\leq\sqrt{106\log2}\), whereas a standard normal has positive
probability beyond every finite bound by its strictly positive density.
Thus the finite construction cannot literally have that continuous law. The theorem
justifies the intended simulation algorithm; it is separate from
implementation replay and range proofs.

\section{Estimators and the complete native results}
For the computed standardized samples \(Z_1,\ldots,Z_T\), the native
driver reports
\[
 \overline Z=\frac{S_1}{T},\qquad
 S_Z^2=\frac{S_2-S_1\overline Z}{T-1},\qquad
 S_1=\sum_j Z_j,\quad S_2=\sum_j Z_j^2,
\]
and \(K_b/T\), where \(K_b=\sum_j\mathbf1_{Z_j>b}\) for
\(b=2,3,q\). The variance numerator equals
\(\sum_j(Z_j-\overline Z)^2\), by expanding the square and using
\(\sum_jZ_j=T\overline Z\). Thus the estimator is the same unbiased
\(T-1\) sample variance used in the Dask pool. Unbiasedness and the
Gaussian mean/chi-square reference laws are proved in
Proposition~\ref{sm:normal-sample}; event moments are proved in
Proposition~\ref{sm:binomial}. Floating-point accumulation and printed
rounding remain implementation approximations to these exact formulas.

\begingroup\small

\begingroup\small
\setlength{\tabcolsep}{4pt}
\Needspace{6\baselineskip}
\noindent\textbf{Retained Ada run: one million observations per size}
\begin{longtable}{rrrrrr}
\toprule
$n$ & Archived mean & Archived variance & $K_2$ & $K_3$ & $K_q$ \\
\midrule
\endfirsthead
$n$ & Archived mean & Archived variance & $K_2$ & $K_3$ & $K_q$ \\
\midrule
\endhead
\bottomrule
\endfoot
$1$ & $3.969606\times10^{-4}$ & $0.999023$ & $22603$ & $1312$ & $1$ \\
$2$ & $5.006006\times10^{-4}$ & $0.999649$ & $22704$ & $1331$ & $0$ \\
$3$ & $5.491356\times10^{-4}$ & $0.999999$ & $22734$ & $1337$ & $1$ \\
$4$ & $5.780900\times10^{-4}$ & $1.000226$ & $22721$ & $1326$ & $1$ \\
$5$ & $5.976301\times10^{-4}$ & $1.000387$ & $22742$ & $1323$ & $1$ \\
$10$ & $6.444156\times10^{-4}$ & $1.000799$ & $22752$ & $1306$ & $1$ \\
$100$ & $7.115883\times10^{-4}$ & $1.001462$ & $22734$ & $1322$ & $1$ \\
$1000$ & $7.290209\times10^{-4}$ & $1.001649$ & $22685$ & $1325$ & $1$ \\
\end{longtable}
\endgroup

\endgroup
These eight rows are the full native simulation output, not selected
examples. Means range from approximately \(0.000397\) to \(0.000729\),
against mean target zero and ideal SE \(0.001\). Composite variances range
from \(0.999023\) to \(1.001649\), against target one. Their largest
absolute variance deviation is about \(0.001649\); the ideal variance
SE is approximately \(0.001414\). For completeness, the variance-SE
formula follows without an additional sampling assumption beyond
Proposition~\ref{sm:normal-sample}: write its chi-square variable as a sum
of \(T-1\) independent squared standard normals. For a standard normal
\(H\), integration by parts using \(\phi'(x)=-x\phi(x)\) gives
\(\E H^4=3\E H^2=3\); the boundary term
\(x^3\phi(x)\) vanishes at both infinities. Hence
\(\Var(H^2)=3-1=2\), and independence gives chi-square variance
\(2(T-1)\). Scaling by \((T-1)^{-2}\) yields
\(\Var(S_Z^2)=2/(T-1)\), whose square root is the stated SE.
The independent-normal representation is the referenced Gaussian
orthogonal-transformation proof in Proposition~\ref{sm:normal-sample}.
These finite-run deviations are retained
as observed, not rounded back to their theoretical values.

Moderate-tail counts range from \(22603\) to \(22752\) for \(Z>2\),
and from \(1306\) to \(1337\) for \(Z>3\). The respective expected
counts are approximately \(22750.132\) and \(1349.898\), obtained by
multiplying the targets in \eqref{eq:hf-tail-targets} by one million.
The native validator records 95\% Wilson intervals for all 16
moderate-tail proportions; all contain their corresponding Gaussian
reference. For example, the \(n=1\), \(Z>2\) interval is approximately
\([0.0223135,0.0228962]\), and its \(Z>3\) interval is
\([0.00124294,0.00138489]\). The score-inversion derivation and
approximate-coverage qualification are in Proposition~\ref{sm:wilson}.

The rare count is zero at \(n=2\) and one at every other size, against
expected count one for each fixed size. By the binomial variance formula,
its ideal count SD is \(\sqrt{1-10^{-6}}\), approximately one, and the
exact probability of no exceedances is \((1-10^{-6})^{10^6}\),
approximately \(e^{-1}\). Proposition~\ref{hfpr:simulation} proves the
zero-count identity and bounds the exponential approximation. The eight
coupled rows must not be pooled as eight million independent tail draws.

\section{Deterministic calculations accompanying the simulation}
\label{sec:hf-spark-calculations}
The companion \path{calculations.csv} supplies nine threshold rows:
the eight finite sizes and the limit \(n\to\infty\). Their reported
variance factors run from \(1\) at \(n=1\), through \(0.64\) at
\(n=10\), to \(0.6004\) at \(n=1000\), with limit \(0.6\).
The corresponding rare raw-mean cutoffs are approximately
\(4.753424309\), \(3.802739447\), \(3.683213761\), and
\(3.681986637\). These are calculations of \(q\sqrt{v_n}\), not
additional Monte Carlo estimates. Proposition~\ref{hfpr:thresholds}
proves the cutoff identity and limit.

Four further outputs check the rare quantile and worked examples. The
three-test mean of 160 on an SD-16 reporting scale gives composite value
\(170.064904975\). The latent conditional example reports mean
\(168.794145026\) and SD \(6.822422923\), on the separately specified
latent SD-16 scale. Proposition~\ref{hfpr:numerics} derives the exact
formulas and numerical interpretation. The normal and optimized native
validators report maximum absolute numerical error
\(2.842170943040401\times10^{-14}\), against tolerance
\(5\times10^{-11}\), over the nine rows and four outputs. This is
agreement with the validator's computed references; it is not a certified
last-digit bound on the external numerical libraries.

\section{Runtime replay, native proof and their boundaries}
The native runtime record reports successful GNAT 13.2.0 builds with
ordinary assertion and overflow checks enabled, both with Ghost Check
and Ghost Ignore. The two policies reproduce the same calculation and
simulation CSV bytes. Normal and optimized Python validators and the
separate-implementation replay pass; replay recovers all three tail
counts exactly in every size and matches means and unbiased variances
within its \(2\times10^{-12}\) tolerance. Eight CLI/domain cases are
also retained for each policy, including rejection of trial counts zero
and one. A repeated policy run or deterministic replay does not add new
random observations.

For the source hashes in that campaign, the native Hoeflin portfolio
discharged 404 checks: 129 by flow analysis and 275 by provers, with zero
unproved and zero justified checks. Native coverage records 53
\texttt{all} and 14 \texttt{spec} entities. The historical structural
inventory covers 35 functions and 28 procedures with 356 explicit
contract entries and a five-entry floor. These are different inventories:
entity/body classification, structural contract counting and discharged
proof obligations are not interchangeable measures. Nor does the
five-entry floor establish logical independence or adequacy of contracts.

The proved scope includes specified modular state advancement, bounded
arithmetic/accumulations, count updates, control flow and modeled report
receipts. Trusted numerical completions, \texttt{erfc}, logarithm,
square root, sine/cosine, compiler/runtime behavior and physical I/O retain
the boundaries documented in the campaign. Receipt contracts establish
modeled requests or normal returns, not successful physical output. No
proof here establishes generator independence, exact normal sampling,
inverse-tail accuracy or validity of the psychometric model.

The present source tree has since moved shared numerical/sampling code
into \path{library/spark/}. The 4 November counts remain evidence about
their pinned source snapshot; they are not a GNATprove result for the
refactored bodies. This report revision reuses immutable evidence and
adds exposition and comparisons. Current-source proof requires a new
matching campaign rather than silently substituting historical Ada.
